\documentclass[12pt, a4paper, UTF8, openany]{ctexbook}

\usepackage{geometry}%设置页边距
\geometry{left=1.54cm, right=1.54cm, top=2.18cm, bottom=2.18cm}

\usepackage{tikz}%绘图
\usetikzlibrary{calc}%绘图时计算
\usetikzlibrary{angles}%角度标识
\usetikzlibrary{intersections}%计算交点

\usepackage{graphicx}%增强图像支持
\usepackage[svgnames]{xcolor}

\usepackage{amsmath, amssymb}%数学

\usepackage{array}%表格基础设置
\usepackage{longtable}%跨页表格
\usepackage{booktabs}%表格水平线

\usepackage{multicol}%多列

\usepackage[colorlinks = true, 
            linkcolor = black]{hyperref}%超链接跳转

\newcommand{\dif}{\mathrm{d}}
\newcommand{\dift}{\,\dif t}

\newcommand{\req}{R_{\text{eq}}}
\newcommand{\eq}[1]{#1_{\text{eq}}}

\newcommand{\unit}{\mathrm{j}}
\newcommand{\re}{\mathfrak{Re}}
\newcommand{\im}{\mathfrak{Im}}

\linespread{1.05}%行间距
\setlength{\parskip}{8pt}%段落间距
\setlength{\parindent}{0pt}%不缩进

\setcounter{chapter}{-1}

\title{Note for Circuit}
\author{1001}

\begin{document}
    \maketitle

    \tableofcontents % use LaTeXmk
    \cleardoublepage % 正文从新页开始，利好双面打印

    \chapter{基础}
    \section{对偶}
    同一侧性质类似，同时把一侧对偶到另一侧。
    \begin{longtable}[h]{ccc}
        并联 & $\leftrightarrow$ & 串联 \\
        电导 & $\leftrightarrow$ & 电阻 \\
        电容 & $\leftrightarrow$ & 电感 \\
        阻抗 & $\leftrightarrow$ & 导纳
    \end{longtable}

    \chapter{基尔霍夫定律}

    \chapter{基本分析方法}

    \chapter{动态电路}
    \section{动态元件}
    这里限制 $u, i$ 均有界，因此积分上限函数总连续，即不能突变。

        \subsection{电容}
        $$\begin{aligned}
            q = Cu & \implies \boxed{i(t) = C \frac{\dif u}{\dif t}} \\
            & \implies u(t) = \frac{1}{C} \int_{t_0}^{t} i \dift
        \end{aligned}$$

        {\heiti 电容电压}不能突变，$u(t^-) = u(t^+)$

        $p = ui = u \cdot C \frac{\dif u}{\dif t} \; \implies \; W = \frac{1}{2} Cu^2|_{t_0}^{t}$

        串联：$u = \sum u_i \; \implies \; \frac{1}{C} \int i \dift = \sum \frac{1}{C_i} \int i \dift \quad \implies \quad \frac{1}{C} = \sum\frac{1}{C_i}$

        并联：$i = \sum i \; \implies \; C \frac{\dif u}{\dif t} = \sum C_i \frac{\dif u}{\dif t} \quad \implies \quad C = \sum C_i$

        \subsection{电感}
        $$\begin{aligned}
            \Psi = Li & \implies \boxed{u(t) = L \frac{\dif i}{\dif t}} \\
            & \implies i(t) = \frac{1}{L} \int_{t_0}^{t} u \dift
        \end{aligned}$$

        {\heiti 电感电流}不能突变，$i(t^-) = i(t^+)$

        $p = ui = L \frac{\dif i}{\dif t} \cdot i \; \implies \; W = \frac{1}{2} Li^2|_{t_0}^{t}$

        串联：$u = \sum u \; \implies \; L \frac{\dif i}{\dif t} = \sum L_i \frac{\dif i}{\dif t} \quad \implies \quad L = \sum L_i$

        并联：$i = \sum i \; \implies \; \frac{1}{L} \int u \dift = \sum \frac{1}{L_i} \int u \dift \quad \implies \quad \frac{1}{L} = \sum\frac{1}{L_i}$

    \section{换路定则}
    电容电压，电感电流，前后均不变。其它均可突变。

    求换路前的 $u_c, i_L$，将其作为初值求换路后。它们初值已知，则此后电路状态均可定，故称为状态变量。

    \section{一阶动态电路}
    一阶：用一个一阶微分方程描述。RC 或 RL。

    电路仍具有线性性，故依叠加定理可将全相应拆分为零输入和零状态的代数叠加。

    LC 电路为一阶微分方程组，至少二阶。

        \subsection{零输入}
        电路的能量仅源自动态元件初始储能，无电压源/电流源。

        以状态变量列微分方程，形如 $x + \tau\frac{\dif x}{\dif t} = 0$。
        特征根 $\lambda = -\frac{1}{\tau}$，通解 $x = x_{0^+}\exp(\lambda t) = x_{0^+}\exp(-\frac{t}{\tau})$。
        
        $\tau$ 具有和 $t$ 相同的量纲，称时间常数。

        \begin{itemize}
            \item RC 电路：$\tau = RC$
            \item RL 电路：$\tau = \frac{L}{R}$
        \end{itemize}
        由对偶，$C \sim G$，$L \sim R$。

        电路中{\heiti 所有量}均以指数变化，故都有通解 $$x = x_{0^+}\exp(-\frac{t}{\tau})$$

        \subsection{零状态}
        仅电压源/电流源供能，动态元件无初始储能。

        以状态变量列微分方程，形如 $k + \tau\frac{\dif x}{\dif t} = x$，其中 $k$ 为常数（电压源/电流源）。
        
        特征根 $\lambda = -\frac{1}{\tau}$，零解 $x = \text{Const} \cdot \exp(-\frac{t}{\tau})$。时间常数与零输入相同。

        只需求任意一个特解，显然 $x \equiv k$ 为一特解，故通解 $x = k + \text{Const} \cdot \exp(-\frac{t}{\tau})$。在 $0^+$ 时刻 $x = 0$，则 $\text{Const} = -k$。进一步地，$t \to \infty, \, x \to k$，于是{\heiti 仅状态变量}有通解 $$x = x_{\infty}(1 - \exp(-\frac{t}{\tau}))$$

        相当于给电容/电感充电，电能转化为磁场能，效率均 $50 \%$。

        \subsection{全相应}
        全相应为零输入和零状态的代数叠加。

        对状态变量，$x = x_{0^+}\exp(-\frac{t}{\tau}) + x_{\infty}(1 - \exp(-\frac{t}{\tau})) = x_{\infty} + (x_{0^+} - x_{\infty})\exp(-\frac{t}{\tau})$

        对任意变量，也必然遵循同样的指数衰减规律。整个电路共用同一个时间常数 $\tau$。故只需求出 $x_{0^+}$ 与 $x_{\infty}$。

        $$x = x_{\infty} + (x_{0^+} - x_{\infty})\exp(-\frac{t}{\tau})$$

        也可看作稳态响应与暂态响应的代数叠加。

        对于一般电路，需先对除动态元件外的电路等效转化，用 $\req$ 计算 $\tau$。

    \chapter{正弦稳态电路}
        \section{基础}
        形式 $f(t) = F_m \cos(\omega t + \Psi)$，电气中统一用{\heiti 余弦}。$F_m$ 振幅，$\omega$ 角频率，$\Psi$ 初相角/初相。

        相位差 $\varphi_{12} = \Psi_1 - \Psi_2$，显然仅对同频变量有定义。$\varphi_{12} >/< 0$ 称 $1$ 超前/滞后 $2$。$\varphi_{12} = 0$ 称同向，$\varphi_{12} = \pm\pi$ 称倒相/反相。

        有效值是对{\heiti 做功}的等效，$F = \sqrt{\frac{1}{T} \int_{0}^{T} f^2(t)\dift} = \frac{1}{\sqrt{2}} F_m$。交流电表示数为有效值，220V 也为有效值。

        符号：
        \begin{longtable}[h]{c | c}
            瞬时值 & $u, i$ \\
            \midrule
            最大值 & $U_m, I_m$ / $U_M, I_M$ \\
            \midrule
            有效值 & $U, I$ \\
            \midrule
            相量 & $\dot{U}, \dot{I}, \dot{U_m}, \dot{I_m}$
        \end{longtable}

        虚数单位 $\unit$，因为 $\mathrm{i}$ 易有歧义。$z = a + b\unit = |z|\angle\theta$。
        
        \section{相量表示}
        欧拉公式：$\exp(\unit\theta) = \cos\theta + \unit\sin\theta$

        定义 $f(t) = F_m \cos(\omega t + \Psi)$ 的相量 $$\dot{F_m} := F_m\angle\Psi = F_m\cos(\omega t + \Psi) + F_m\sin(\omega t + \Psi)\unit$$

        进一步地，有效值相量 $$\dot{F} := \frac{1}{\sqrt{2}} \dot{F_m} = F\angle\Psi$$

        相量的单位与原单位相同，$f(t) = \re[\dot{F_m} \cdot \exp(\unit \omega t)]$。若整个电路均同频，则所有相量均以同样的角速度旋转，相对位置并不改变，因此无需考虑时间。$\exp(\unit \omega t)$ 称旋转因子。

        \subsection{基尔霍夫定律}
        KCL: $\sum \dot{I} = 0$

        KVL: $\sum \dot{U} = 0$

        证明：$\sum \re[ \dot{F}_i \exp(\unit \omega t) ] = \re[ \sum \dot{F}_i \exp(\unit \omega t) ] = \re[ (\sum \dot{F}_i) \exp(\unit \omega t) ] = 0 \implies \sum \dot{F}_i = 0$

        \subsection{电阻、电容、电感}
        \begin{itemize}
            \item 电阻 \\
                $\dot{U} = \dot{I} R, \; \dot{I} = \dot{U} G$
        \end{itemize}

        $f = F_m \cos(\omega t + \Psi)$

        $g = k\frac{\dif f}{\dif t} = k\omega F_m \cos(\omega t + \Psi + \frac{\pi}{2}) \; \implies \dot{G} = \omega k |\dot{F}|\angle (\Psi + \frac{\pi}{2}) = \unit \omega k \dot{F}$

        \begin{itemize}
            \item 电容 \\
                定义容抗 $X_C := \frac{1}{\omega C}$，单位欧姆 $\Omega$。
                $$\frac{\dot{U}}{\dot{I}} = \frac{1}{\unit \omega C} = \frac{1}{\omega C} \angle -\frac{\pi}{2} = -\unit X_C$$
                电流超前电压 $90\%$，阻低频，同高频。
            \item 电感 \\
                定义感抗 $X_L := \omega L$，单位欧姆 $\Omega$。
                $$\frac{\dot{U}}{\dot{I}} = \unit \omega L = \omega L \angle \frac{\pi}{2} = \unit X_L$$
                电流滞后电压 $90\%$，阻高频，同低频。
        \end{itemize}
        
        一般地，定义：\\
        阻抗 $Z = \frac{\dot{U}}{\dot{I}} = R + \unit X$，单位欧姆 $\Omega$ \\
        导纳 $Y = \frac{\dot{I}}{\dot{U}} = G + \unit B$，单位西门子 S

        $Z, Y$ 均为复数，阻抗角 $\varphi_Z = \arg Z$ 与导纳角 $\varphi_Y = \arg Y$ 互为相反数。

        对任意网络，可等效为 $R$ 和 $\unit X$ 串联，或 $G$ 和 $\unit B$ 并联。\\
        {\heiti 倒数关系不成立}：$R \not\equiv \frac{1}{G}$，$X \not\equiv \frac{1}{B}$

        无源网络 $\varphi_Z$: \\
        a. $0$：纯电阻 \\
        b. $(-\frac{\pi}{2}, 0)$ / $-\frac{\pi}{2}$：容性 / 纯电容性 \\
        c. $(0, \frac{\pi}{2})$ / $\frac{\pi}{2}$：感性 / 纯电感性 \\
        有源网络 $|\varphi_Z|$ 可以超过 $\frac{\pi}{2}$

        电容、电感当复数电阻计算。

        \section{功率}
        W = var = V$\cdot$A，三个单位量纲相同，但用于不同功率。

        \subsection{瞬时功率}
        $$\begin{aligned}
            p(t) &= u(t) i(t) \\
            &= U_m I_m \cos(\omega t + \Psi_u) \cos(\omega t + \Psi_i) = U_m I_m \cdot \frac{1}{2}[ \cos(\Psi_u - \Psi_i) + \cos(2 \omega t + \Psi_u + \Psi_i) ] \\
            &= \boxed{UI \cos\varphi_{ui}} + \boxed{UI\cos(2 \omega t + \Psi_u + \Psi_i)} \\
            &= UI \cos\varphi_{ui} + UI \cos(2 \omega t + 2\Psi_u - \varphi_{ui}) \\
            &= UI \cos\varphi_{ui} + UI [\cos 2(\omega t + \Psi_u) \cos \varphi_{ui} + \sin 2(\omega t + \Psi_u) \sin\varphi_{ui}] \\
            &= \boxed{UI \cos\varphi_{ui} [1 + \cos 2(\omega t + \Psi_u)]} + \boxed{UI \sin\varphi_{ui} \sin 2(\omega t + \Psi_u)}
        \end{aligned}$$

        $u, i$ 相位差 $\varphi_{ui}$

        基准 $UI \cos\varphi_{ui}$，偏移 $UI \cos(2 \omega t + \Psi_u + \Psi_i)$，频率为电压/电流频率之二倍。

        有功分量 $UI \cos\varphi_{ui} [1 + \cos 2(\omega t + \Psi_u)]$，恒非负；\\
        无功分量 $UI \sin\varphi_{ui} \sin 2(\omega t + \Psi_u)$，在一个周期内总共归零。

        \subsection{平均功率与无功功率}
        三角函数在一个周期内积分必为零。

            \subsubsection{平均功率 / 有功功率}
            $$P = UI \cos\varphi_{ui}$$

            单位瓦特 W，描述电阻的总功率。

            a. $Z = \eq{R} + \unit\eq{X} \; \implies \; P = I^2 \eq{R}$ \\
            b. $Y = \eq{G} + \unit\eq{B} \; \implies \; P = U^2 \eq{G}$

            \subsubsection{无功功率}
            $$Q = UI \sin\varphi_{ui}$$

            单位乏 var，描述电容、电阻的总功率，也即能量交换的最大功率。

            a. $Z = \eq{R} + \unit\eq{X} \; \implies \; Q = I^2 \eq{X}$ \\
            b. $Y = \eq{G} + \unit\eq{B} \; \implies \; Q = -U^2 \eq{B}$

            由 $\sin\varphi_{ui}$ 可见，容性支路 $Q < 0$，感性支路 $Q > 0$。

        \subsection{视在功率与功率因数}
        视在功率 $$S = UI$$

        单位伏安 V$\cdot$A，描述理论最大功率，也即电器的额定功率。

        功率因数 $$\frac{P}{S} = \cos\varphi_{ui}$$

        一般地，电路多呈感性，需配电容以减小 $\varphi_{ui}$，进而增大功率。

        \subsection{最大功率传输}

        等效电压源 $U$，等效内阻抗 $Z = R + \unit X$，则外电路等效阻抗 $Z' := \bar{Z} = R - \unit X$ 时功率最大。证明只需求导。

        $P_{\max} = \frac{U^2}{4R}$

        \section{谐振}
        $u, i$ 同相，即 $Z/Y$ 无虚部。

        $$\omega_0 = \frac{1}{\sqrt{LC}}$$

\end{document}